\newpage
\setcounter{table}{0}
\setcounter{figure}{0}
\section{总结与展望}
\subsection{研究总结}
本研究针对奖励型众筹项目成功性预测这一核心问题，突破了传统研究仅关注内容语义的局限性，创新性地引入了内容呈现结构这一维度，构建了基于统一序列建模范式的多模态深度学习模型，并在 Kickstarter 真实数据集上进行了实验验证。现将本研究的核心内容总结如下：

首先，本研究立足叙事传输理论，设计了基于统一序列建模范式的多模态深度学习模型 MSTC，实现了对众筹界面内容呈现结构这一维度的量化建模。该模型摒弃了传统研究中将模态视为孤立集合的局限，利用 HTML DOM 树解析技术还原页面真实的物理展示顺序。在底层表征上，模型中的 Token 编码器融合了 CLIP 语义嵌入、区分不同模态的类型向量和表征认知负荷的视觉密度属性，实现了对内容块的多维信息封装。通过正弦位置编码的注入，Token 序列被赋予了感知真实页面浏览顺序的能力，进而使得 Transformer 序列编码器可以利用多头自注意力机制模拟读者认知行为，捕捉图文交织中所蕴含的叙事逻辑与视觉节奏。为了提升模型在众筹数据高噪声、分布不均等复杂场景下的性能，本研究系统性地应用了标签平滑、指数滑动平均（EMA）以及动态阈值寻优等一系列稳健性策略，显著增强了预测结果的泛化能力与评估稳定性。

其次，基于 Kickstarter 真实数据集的实验结果证明，统一序列建模范式在众筹结果预测任务上具有显著优越性。在与多种主流多模态融合范式的对比中，本研究提出的 MSTC 模型在准确率、F1 分数以及 AUC 等核心指标上均表现最优，且在精确率与召回率之间取得了最佳平衡。这种性能增益源于模型 Token 级的深度跨模态交互，该机制允许模型在特征提取早期捕获图文间的深层对齐关系，最大限度地减少了特征压缩过程中的语义损失。消融实验结果进一步支持了本研究的核心假设，即图文内容呈现结构是影响众筹成败的关键判别变量，众筹页面吸引潜在投资者的能力不仅源于发起人“说了什么”，更取决于发起人“是如何说的”。消融实验显示，引入正弦位置编码并遵循众筹页面中真实的图文展示顺序，能为模型带来显著且可复现的性能增益；视觉密度特征是语义信息之外的重要补充，能够稳定提升模型对复杂页面信息的判别能力。

综上所述，本研究不仅提升了深度学习模型进行众筹成功性预测的性能，更从实证角度揭示了图文内容呈现结构在奖励式众筹这一金融模式中的重要地位，为平台优化资源配置和发起人改进项目设计提供了科学的决策依据。

\subsection{理论贡献}
本研究对众筹页面多模态内容呈现结构进行建模，不仅在技术层面提升了预测性能，更在理论层面深化了对众筹交互机制的理解。

这种理解的深化，首先源于对众筹内容呈现结构信号作用的挖掘。不同于以往研究多关注发起人背景及图文内容质量，本研究通过实验发现，页面内容的呈现形式本身也是反映众筹项目质量的关键判别信号。在奖励型众筹这种典型的信息不对称场景下，一套有序、严谨且符合逻辑的内容结构，反映了发起人在项目策划上的投入程度与专业性，进而影响了潜在投资者的出资意愿。这一发现拓展了众筹结果预测领域特征选择的边界，打破了过往研究只关注内容而不关注结构的局限性，为众筹项目质量评估提供了新视角。

而在探究结构如何发挥作用的过程中，本研究为叙事传输理论提供了计算层面的工程实证方案，实现了从理论驱动到工程实践的闭环。这种闭环的建立源于对叙事传输质量在计算维度的建模：高质量的众筹界面不仅在顺序上需要有严谨的递进逻辑，同时在密度上需要保持具有呼吸感的叙述节奏。通过构建深度学习模型 MSTC，本研究量化了图文序列的叙事逻辑和视觉密度，通过消融实验证明了遵循真实展示顺序的序列建模能显著提升预测性能、视觉密度特征也能通过量化认知负荷实现稳健的判别增益。这一结果在计算层面有力地支撑了叙事传输机理在信息传播与说服过程中的核心地位。

本研究模型的有效性，进一步验证了图文统一建模在模拟人类认知过程中的独特价值。基于双重编码理论，文字语义与视觉呈现之间存在深层的协同关系。本研究在实证层面证实了将文字与图片纳入统一序列进行特征融合的必要性，认为该方案比传统的模态割裂建模更符合人类处理交替信息的认知规律。对比实验表明，通过使用 Transformer 编码器在 Token 级别实现图文的深度交互，模型在预测性能上超越了晚期融合模型和卷积神经网络模型，强调了图文统一建模在复杂内容理解中的作用。

\subsection{局限性与未来方向}
虽然本研究成功在众筹成功性预测中引入了内容呈现结构，验证了该方案的有效性，但由于时间、数据及计算资源的限制，仍存在以下局限，同时也为未来的深入研究指明了方向。

第一，未来研究可以涵盖视频模态。当前研究主要聚焦于图文交织的界面布局结构，未考虑视频模态，主要原因在于：在实际的众筹决策场景中，长达数分钟的项目视频虽然能提供丰富的感官信息，但其信息获取成本远高于图文阅读。大量投资者倾向于通过快速滑动页面来获取项目的关键信息，而非完整观看视频。然而不可忽视的是，视频往往位于项目封面等核心位置，其动态演示具有极强的感染力，现已有部分研究将视频模态引入深度学习预测模型并取得了良好成效。未来的研究中可以尝试提取视频关键帧进行语义理解和情感分析，将其作为特殊的视觉 Token 整合进统一序列框架中。

第二，跨平台泛化性有待验证。本研究的实验数据集仅来源于 Kickstarter 平台，虽然其作为全球最大的奖励型众筹平台具有代表性，但在平台规则、用户画像及审美倾向方面可能存在特定偏好。考虑到不同平台的叙事风格可能存在差异，不同国家的众筹平台也可能存在不同的用户群体，未来可以将模型迁移至其他国内外众筹平台（如 Indiegogo 等），验证基于内容呈现逻辑的模型在不同平台上的泛化能力，提升模型的通用性。

第三，多模态特征提取的精细度仍可提升。目前模型主要依赖 CLIP 模型提取的语义特征和基于物理属性的视觉密度，这种建模方式虽然高效，但存在特征提取不够细腻的问题。首先，众筹场景下的图文语义分布可能和 CLIP 预训练时所使用的图像及文本存在差异，未来研究可尝试使用众筹领域特有的图文数据对 CLIP 进行微调，使模型能够学习到众筹场景下的专业词汇与视觉逻辑，从而提取更具判别力的语义特征。其次，除图文语义信息之外，图文的质量属性也是叙事质量的重要组成部分。未来可引入专业的美学评估模型量化图像的视觉吸引力指标，也可利用大语言模型对文本质量（如流畅度、可读性等）进行深度评估，更全面地刻画众筹项目的叙事质量。
